\section{椭圆曲线}\label{sec:8.4}
\setcounter{equation}{11}
\setcounter{figure}{23}

\SourceParagraph{ch08-8.4-intro-001}{626}
现在证明椭圆曲线 HMS 的一个受限版本。为简单起见，不尝试比较完整的导出范畴与 Fukaya 范畴，不过，这里的工作已包含其中大部分。椭圆曲线情形已在\BookCite{ref397,ref396}中详细研究；这里采用稍有不同的方法，受到 Abouzaid\BookCite{ref002}启发，并在\BookCite{ref001}中发展，利用 Čech 方法定义导出范畴的 $A_\infty$ 结构。

\SourceParagraph{ch08-8.4-intro-002}{626}
将构造椭圆曲线的退化族 $\mathcal X\to D$。$D$ 可以看作坐标为 $q$ 的单位圆盘；在概形范畴中，则取 $D=\operatorname{Spec}\mathbb C[[q]]$。此族为\BookSectionRef{sec:6.2.4}{6.2.4}中所考虑族的紧化。使用这样的族的优点，是参数 $q$ 对应于镜像 Fukaya 范畴中的 Kähler 参数 $q$，也就是 Novikov 环的变量。特别地，应该将两侧都看作由 $q$ 参数化的范畴族。此外，采用显式退化，使这里的方法与 Gross、Siebert\BookCite{ref205,ref206}提出的思想很好地契合。虽然不在这里展开这些思想，它们在这一情形仍很有帮助。

\SourceParagraph{ch08-8.4-intro-003}{626}
实际上，将在构造退化 $\mathcal X\to D$ 的同时，构造线丛 $\mathcal L$，其在每条纤维上的限制都丰沛，即正。随后，计算涉及 $\mathcal L$ 的正、负张量幂的全部高阶乘法映射。

\subsection{退化：Tate 曲线}\label{sec:8.4.1}

\SourceParagraph{ch08-8.4.1-001}{627}
所用退化构造方法，是 Mumford 关于 Abel 簇退化的构造\BookCite{ref373}的一个特例。这一椭圆曲线退化通常称为 Tate 曲线。从下述数据出发：选择线丛的正整数次数 $d$，考虑连续分段线性函数 $\varphi\colon\mathbb R\to\mathbb R$，在每个区间 $[i,i+1]$ 上斜率为 $i$，$i\in\mathbb Z$。具体地，可以令它在 $[0,1]$ 上取值零；此时，
\[
\varphi(x)=ix-\frac{i(i+1)}2,\qquad x\in[i,i+1],
\]
且满足周期性条件
\begin{equation}\label{eq:8.12}
\varphi(x+d)=\varphi(x)+dx+\frac{d(d-1)}2.
\end{equation}
令 $\Delta=\{(x,y)\in\mathbb R^2\mid y\geq\varphi(x)\}$。
\begin{figure}[H]
\centering
\BookImage{01_Numbered_Figures/Chapter_08/Ch08_Fig24_PDFp627_Bookp617.pdf}
\caption{$\varphi$ 的图像。阴影区域为 $\Delta$。}\label{fig:8.24}
\end{figure}

\SourceParagraph{ch08-8.4.1-002}{627}
正如凸多胞体决定射影环面簇，无界多面体也可以决定环面簇。这可通过 $\Delta$ 的法扇描述。对于每个面 $\sigma\subseteq\Delta$，定义 $\Delta$ 在 $\sigma$ 处的法锥
\[
\check\sigma=\left\{n\in(\mathbb R^2)^\vee\ \middle|\
\begin{gathered}
n|_\sigma\text{ 为常数},\\
\langle n,m\rangle\geq\langle n,m'\rangle
\quad\forall m\in\Delta,\ m'\in\sigma
\end{gathered}\right\}.
\]
法扇为各面对应的锥的集合 $\{\check\sigma\mid\sigma\subseteq\Delta\text{ 为面}\}$，位于 $(\mathbb R^2)^\vee$ 中。容易看出，若 $\sigma$ 是端点为 $(i,\varphi(i))$、$(i+1,\varphi(i+1))$ 的区间，则 $\check\sigma$ 为 $(-i,1)$ 生成的射线。因此，$\Delta$ 的法扇 $\Sigma$ 如\BookRef{fig:8.25}{图25}所示。它包含无穷多个锥，超出了通常对扇的有限性要求，不过，这并不妨碍构造：它仍定义环面簇 $X_\Sigma$，由这些锥给出的无穷多个仿射环面坐标片覆盖。

\SourceParagraph{ch08-8.4.1-003}{628}
此外，$\mathbb R^2$ 中的 $(0,1)$ 在 $\Sigma$ 每条射线的本原生成元上取值一。将对应于 $(0,1)$ 的单项式记为 $q$，则它为 $X_\Sigma$ 上的正则函数，定义映射 $g\colon X_\Sigma\to\mathbb A^1$。
\begin{figure}[H]
\centering
\BookImage{01_Numbered_Figures/Chapter_08/Ch08_Fig25_PDFp628_Bookp618.pdf}
\caption{$\Delta$ 的法扇。}\label{fig:8.25}
\end{figure}

\SourceParagraph{ch08-8.4.1-004}{628}
由于 $q$ 在 $X_\Sigma$ 的每个环面除子上，即每条射线对应的除子上，以一阶消失，纤维 $g^{-1}(0)$ 为无穷多个 $\mathbb P^1$ 组成的链。另一方面，不难看出，$g$ 的其他每条纤维，均为 $\mathbb C^*$ 的副本。

\SourceParagraph{ch08-8.4.1-005}{628}
像任意环面簇一样，$X_\Sigma$ 有由扇的锥，等价地由 $\Delta$ 的面标记的开覆盖。特别地，有开覆盖 $\{U_w\mid w\text{ 为 }\Delta\text{ 的顶点}\}$，其中
\[
U_w=\operatorname{Spec}\mathbb C[\check w^\vee\cap\mathbb Z^2].
\]
它是对应于扇中锥 $\check w$ 的开子集。也可以将 $\check w^\vee$ 看作 $\Delta$ 在 $w$ 处的切楔，见\BookRef{fig:8.26}{图26}。同样，若 $\sigma$ 为 $\Delta$ 的边，则得到开子集 $U_\sigma=\operatorname{Spec}\mathbb C[\check\sigma^\vee\cap\mathbb Z^2]$。
\begin{figure}[H]
\centering
\BookImage{01_Numbered_Figures/Chapter_08/Ch08_Fig26_PDFp628_Bookp618.pdf}
\BookNumberOnlyCaption\label{fig:8.26}
\end{figure}

\SourceParagraph{ch08-8.4.1-006}{628}
现在，$\mathbb Z$ 在 $X_\Sigma$ 上作用，可以通过它在扇 $\Sigma$ 上的作用描述。作用生成元在 $(\mathbb R^2)^\vee$ 上为矩阵
\[
\begin{pmatrix}1&d\\0&1\end{pmatrix},
\]
将射线 $\mathbb R_{\geq0}(i,1)$ 映为 $\mathbb R_{\geq0}(i+d,1)$。在大环面轨道 $(\mathbb C^*)^2\subseteq X_\Sigma$ 上，取对应于 $\mathbb R^2$ 标准基 $(1,0),(0,1)$ 的坐标 $z,q$，它给出作用
\[
(z,q)\longmapsto(zq^d,q).
\]

\SourceParagraph{ch08-8.4.1-007}{629}
因此，固定非零 $q$ 后，作用在由 $z$ 参数化的 $\mathbb C^*$ 上为 $z\mapsto zq^d$。对于 $0<|q|<1$，此纤维的 $\mathbb Z$ 商为 $\mathbb C^*/q^{d\mathbb Z}$，等同于
\[
\mathbb C/\left\langle1,\frac d{2\pi\mathrm i}\log q\right\rangle.
\]
注意，$\exp(2\pi\mathrm i\,\cdot)$ 将 $\mathbb C/\mathbb Z$ 与 $\mathbb C^*$ 识别；只要 $|q|\neq1$，这一商显然为椭圆曲线。

\SourceParagraph{ch08-8.4.1-008}{629}
另一方面，$q=0$ 时，作用将 $\mathbb P^1$ 的无穷链平移 $d$ 个位置。因此，其 $\mathbb Z$ 商为 $d$ 条有理曲线组成的环，如\BookRef{fig:8.27}{图27}所示。
\begin{figure}[H]
\centering
\BookImage{01_Numbered_Figures/Chapter_08/Ch08_Fig27_PDFp629_Bookp619.pdf}
\BookNumberOnlyCaption\label{fig:8.27}
\end{figure}

\SourceParagraph{ch08-8.4.1-009}{629}
现在取 $\mathcal X=g^{-1}(D)/\mathbb Z$，$D$ 为开单位圆盘。由 $q$ 诱导的 $g\colon\mathcal X\to D$，为椭圆曲线向所谓 $I_d$ 型 Kodaira 纤维的退化，即向 $d$ 条有理曲线组成的环退化。

\SourceParagraph{ch08-8.4.1-010}{629}
这里，就像 Fukaya 范畴中使用作为形式幂级数环的 Novikov 环一样，采用此商的形式版本更方便。这一点技术性较强，但它使我们基本上可以完全围绕奇异纤维的拓扑工作，也使 $\mathbb Z$ 商能够在更适合代数几何的范畴中构造。

\SourceParagraph{ch08-8.4.1-011}{629}
确实，很难对代数簇直接取 $\mathbb Z$ 这类群的商。这里先注意，$g^{-1}(0)/\mathbb Z$ 作为代数簇有意义；如前所述，它为 $d$ 条有理曲线组成的环。另一方面，令 $X_\Sigma^k$ 为方程 $q^{k+1}=0$ 定义的子概形，即 $g^{-1}(0)$ 的“$k$ 阶加厚”。可以验证，$X_\Sigma^k/\mathbb Z$ 也作为概形有意义，为 $g^{-1}(0)/\mathbb Z$ 的 $k$ 阶加厚；具体细节对这里不重要。

\SourceParagraph{ch08-8.4.1-012}{630}
随后构造形式概形 $\widehat X_\Sigma$，即沿 $g^{-1}(0)$ 的形式完备化，记为 $\widehat X_\Sigma=\varprojlim X_\Sigma^k$。形式概形的定义见\BookCite{ref222}第 II 章第 9 节。它是赋环空间，底层拓扑空间为 $X_\Sigma^k$ 的拓扑空间，对所有 $k$ 均相同；等价地，就是 $g^{-1}(0)$。对于开集 $U\subseteq g^{-1}(0)$，其结构层满足
\[
\Gamma(U,\mathcal O_{\widehat X_\Sigma})
=\varprojlim_k\Gamma(U,\mathcal O_{X_\Sigma^k}).
\]
逆极限的定义以及环完备化的更多细节，见例如\BookCite{ref025}。

\SourceParagraph{ch08-8.4.1-013}{630}
本质上，这允许考虑 $g^{-1}(0)$ 的“形式邻域”中的函数，它们是关于 $q$ 的形式幂级数。例如，对于前述开集 $U_w$，$\widehat U_w=g^{-1}(0)\cap U_w$ 在拓扑上是 $g^{-1}(0)$ 的开子集，且
\[
\begin{aligned}
\Gamma(\widehat U_w,\mathcal O_{\widehat X_\Sigma})
&=\varprojlim_n\mathbb C[\check w^\vee\cap\mathbb Z^2]/(q^n)\\
&=: \mathbb C[\check w^\vee\cap\mathbb Z^2]
\widehat\otimes_{\mathbb C[q]}\mathbb C[[q]].
\end{aligned}
\]
右侧的完备张量积由左侧逆极限定义。于是，$\widehat{\mathcal X}:=\widehat X_\Sigma/\mathbb Z$ 也作为形式概形有意义，有由 $q$ 诱导的形式概形映射 $\widehat g\colon\widehat{\mathcal X}\to\widehat{\mathbb A}^1$。这里，$\widehat{\mathbb A}^1$ 为由点零及环 $\mathbb C[[q]]$ 组成的赋环空间，应看作 $g\colon\mathcal X\to D$ 的形式版本。

\SourceParagraph{ch08-8.4.1-014}{630}
实际上，将在 $\widehat{\mathcal X}$ 上计算所需的乘法映射，结果均为关于 $q$ 的形式幂级数。不过，首先需要构造线丛。

\SourceParagraph{ch08-8.4.1-015}{630}
选择 $\widehat{\mathcal X}$ 上的相对丰沛线丛。其实，$\mathcal L$ 的选择已经包含在 $\Delta$ 中：多面体不仅决定环面簇，也决定其上的线丛。这里，$X_\Sigma$ 上有线丛 $\mathcal L=\mathcal O_\Delta(1)$，其截面基与 $\Delta$ 中的整格点一一对应。可以通过环面簇线丛的通常转移映射描述它：$\mathcal L|_{U_w}$ 平凡，自然识别为 $\mathcal O_{U_w}\cdot z^w$。因此，在 $U_w\cap U_{w'}$ 上，从 $\mathcal O_{U_w}$ 到 $\mathcal O_{U_{w'}}$ 的转移映射为乘以 $z^{w-w'}$。容易看出，
\[
\Delta=\bigcap_{w\text{ 为 }\Delta\text{ 顶点}}(\check w^\vee+w),
\]
所以每个整格点 $m\in\Delta$ 都给出 $\mathcal L$ 的良好定义的截面 $z^m$。注意，将 $\mathcal L$ 限制到 $g^{-1}(0)$ 的任一不可约分支，就得到 $\mathcal O_{\mathbb P^1}(1)$。

\SourceParagraph{ch08-8.4.1-016}{630}
希望 $\mathcal L$ 在 $\widehat{\mathcal X}$ 上诱导线丛 $\widehat{\mathcal L}$。为此，仅知道 $\mathbb Z$ 在 $X_\Sigma$ 上的作用还不够，必须知道它如何提升为 $\mathcal L$ 上的作用。只需描述它在
\[
\bigoplus_{n=0}^\infty H^0(X_\Sigma,\mathcal L^{\otimes n})
\]
上的作用。此截面空间有由 $C(\Delta)\subseteq\mathbb R^3$ 的整格点给出的典范基，其中
\[
C(\Delta)=\overline{\{(rm,r)\mid m\in\Delta,\ r>0\}}.
\]
第三坐标为整数 $n$ 的整格点，对应于 $H^0(X_\Sigma,\mathcal L^{\otimes n})$ 的基元素。

\SourceParagraph{ch08-8.4.1-017}{631}
所需作用应满足下列性质。
\begin{enumerate}
\item 在 $\mathbb Z^3\subseteq\mathbb R^3$ 上线性，因为希望它通过 $\mathcal L$ 总空间的环面自同构作用。
\item 保持第三坐标，即不改变截面的次数。
\item 与 $X_\Sigma$ 上的作用相容，因此在 $\mathbb Z^2\times\{0\}\subseteq\mathbb Z^3$ 上的作用，由矩阵 $\begin{pmatrix}1&0\\d&1\end{pmatrix}$ 生成，它是 $(\mathbb R^2)^\vee$ 上作用的转置。
\item 将 $C(\Delta)$ 映到自身。
\end{enumerate}
前三个条件要求生成元为 $\begin{pmatrix}1&0&*\\d&1&*\\0&0&1\end{pmatrix}$。第四个条件进一步要求保持 $C(\Delta)$ 的边界，由周期性\BookEqRef{eq:8.12}{8.12}，唯一的选择为
\[
T=\begin{pmatrix}
1&0&d\\
d&1&\dfrac{d(d-1)}2\\
0&0&1
\end{pmatrix}.
\]

\SourceParagraph{ch08-8.4.1-018}{631}
对于给定 $n\geq0$，将此生成元写为仿射变换
\[
T_n(x,y):=\left(x+nd,xd+y+\frac{d(d-1)}2n\right),
\]
即 $T(x,y,n)$ 的前两坐标。这描述了它在 $\mathcal L^{\otimes n}$ 截面上的作用，截面基由 $n\Delta$ 的整格点代表。

\SourceParagraph{ch08-8.4.1-019}{631}
像先前一样，现对 $\mathcal L$ 取 $\mathbb Z$ 商，这只在 $g^{-1}(0)$ 的形式邻域 $\widehat X_\Sigma$ 中有意义。由此得到 $\widehat{\mathcal X}$ 上的线丛 $\widehat{\mathcal L}$。

\SourceParagraph{ch08-8.4.1-020}{631}
最简单的描述，是写出 $H^0(\widehat{\mathcal X},\widehat{\mathcal L}^{\otimes n})$ 的 $\mathbb C[[q]]$ 截面基，即形式完备化上 $\mathcal L^{\otimes n}$ 的 $\mathbb Z$ 不变截面的基。对于 $n\geq1$，任取 $p\in\tfrac1n\mathbb Z$，定义
\[
\vartheta_{n,p}:=\sum_{s=-\infty}^\infty
z^{T_n^s(np,n\varphi(p))}.
\]
注意，在任意给定开集 $U_w$ 上，此式在形式意义下成立，即
\[
\vartheta_{n,p}\in
\bigl(\mathbb C[\check w^\vee\cap\mathbb Z^2]
\widehat\otimes_{\mathbb C[q]}\mathbb C[[q]]\bigr)\cdot z^{nw}.
\]
对于任意给定 $k>0$，定义此和式的项，除有限个外，在
\[
\bigl(\mathbb C[\check w^\vee\cap\mathbb Z^2]
\otimes_{\mathbb C[q]}\mathbb C[q]/(q^k)\bigr)\cdot z^{nw}
\]
中均为零，由此即可看出其良好定义。这些函数称为 theta 函数。

\SourceParagraph{ch08-8.4.1-021}{632}
还有 $\vartheta_{n,p}=\vartheta_{n,p+d}$，因为
\[
T_n(np,n\varphi(p))=(n(p+d),n\varphi(p+d)).
\]
不难看出，各 $\vartheta_{n,p}$ 作为 $\mathbb C[[q]]$ 模生成 $\mathcal L^{\otimes n}|_{\widehat X_\Sigma}$ 的 $\mathbb Z$ 不变截面空间。不过，必须容许关于 $q$ 的形式幂级数，才能存在这种非零不变截面；对于 $n>0$，$X_\Sigma$ 上的 $\mathcal L^{\otimes n}$ 没有这样的非零截面。这再次说明，为取商，必须先作完备化。更一般地，若 $(x,y)\in\mathbb Z^2\cap n\Delta$，即 $y\geq n\varphi(x/n)$，则
\begin{equation}\label{eq:8.13}
\sum_{s=-\infty}^\infty z^{T_n^s(x,y)}
=q^{y-n\varphi(x/n)}\vartheta_{n,x/n}.
\end{equation}

\SourceParagraph{ch08-8.4.1-022}{632}
定义
\begin{equation}\label{eq:8.14}
\operatorname{ord}_n(x,y)=y-n\varphi(x/n),
\end{equation}
其中 $n\in\mathbb Z\setminus\{0\}$、$(x,y)\in\mathbb Z^2$。若 $n=0$、$x=0$，则自然取 $\operatorname{ord}_0(0,y)=y$。这一概念的动机来自\BookEqRef{eq:8.13}{8.13}。

\SourceParagraph{ch08-8.4.1-023}{632}
这一不变截面基的元素，与整仿射流形 $B=\mathbb R/d\mathbb Z$ 上坐标属于 $\tfrac1n\mathbb Z$ 的点一一对应，见\BookSectionRef{sec:6.1.2}{6.1.2}；马上还将回到这一点。将这些点的集合记为 $B((1/n)\mathbb Z)$。

\SourceParagraph{ch08-8.4.1-024}{632}
若读者不满意上述较粗略的商构造描述，也可以如下处理。考虑 $\mathbb C[[q]]$ 代数
\[
\widehat R=\bigoplus_{n=0}^\infty
H^0(\widehat{\mathcal X},\widehat{\mathcal L}^{\otimes n}).
\]
仅需知道其 $\mathbb C[[q]]$ 模生成元集合
\[
\{1\}\cup\{\vartheta_{n,p}\mid p\in B((1/n)\mathbb Z),\ n\geq1\},
\]
以及乘法映射，也就是代数的结构常数，即可定义它。乘法马上将说明。随后考虑概形 $\operatorname{Proj}\widehat R$，它定义在 $\operatorname{Spec}\mathbb C[[q]]$ 上，参见\BookCite{ref222}第 II 章第 2 节。这是通常意义的概形，而非形式概形。存在丰沛线丛时，可以从形式概形得到通常概形，这正是 Grothendieck 存在性定理，见 EGA III，5.4.5\BookCite{ref211}。所以，上述构造也可看作这一一般结论的特例。

\subsection{theta 函数的乘法与三角形}\label{sec:8.4.2}

\SourceParagraph{ch08-8.4.2-001}{633}
现在转向 theta 函数及其乘法的几何解释。前面对 theta 函数基的描述，已经产生一个仿射流形 $B$，因此，自然应尝试将当前构造放入\BookChapterRef{6}的框架。

\SourceParagraph{ch08-8.4.2-002}{633}
实际上，将 $\mathcal X\to D$ 的纤维看作 $\check X(B)$，见\BookSectionRef{sec:6.2.1}{6.2.1}。纤维带有线丛 $\mathcal L$ 的限制，其次数必为 $d$。另一方面，$\check X(B)$ 有典范辛形式 $\omega$，满足 $\int_{\check X(B)}\omega=d$，因此 $\omega$ 代表 $c_1(\mathcal L)$。

\SourceParagraph{ch08-8.4.2-003}{633}
读者或许认为，这与\BookSectionRef{sec:6.3}{6.3}的思路相反：那里，在 $X(B)$ 上考虑线丛，在 $\check X(B)$ 上考虑 Lagrange 子流形。确实如此，不过这一交换有其原因。讨论 Fukaya 范畴时，至少在椭圆曲线情形，可以使用固定的概复结构，即 $X(B)$ 的复结构；它的全纯圆盘很容易描述。选择在 Fukaya 一侧使用 $X(B)$，使圆盘易于描述，却使 Lagrange 子流形较难描述。同样，在导出范畴一侧使用 $\check X(B)$，层较难描述，却从辛形式获得一个正则给定的上同调类。因此，只考虑第一 Chern 类由 $\omega$ 的整倍数代表的线丛，更具体地，只考虑 $\mathcal L$ 的幂。在 $X(B)$ 上，则只考虑非常特殊的 Lagrange 子流形：由展开映射的倍数诱导的截面，见\BookSectionRef{sec:6.1.2}{6.1.2}，其局部系统平凡。

\SourceParagraph{ch08-8.4.2-004}{633}
按照\BookSectionRef{sec:6.3}{6.3}的思路，更显式地描述与 $\mathcal L^{\otimes n}$ 镜像对偶的 $X(B)$ 上的 Lagrange 子流形。首先，指定 $X(B)$ 的辛形式。由\BookRef{prop:6.14}{命题6.14}，可以通过 $B$ 上的凸多值函数指定；取 $K=y^2/2$ 就很好，得到 $\omega=\mathrm dy\wedge\mathrm dx$。这里暂时沿用第 6 章的约定，以 $y$ 为 $B$ 的坐标，$x$ 为 $X(B)$ 的纤维坐标。任何实曲线都是 Lagrange 的，所以 $\omega$ 在这里唯一的作用，是计算全纯圆盘的面积。

\SourceParagraph{ch08-8.4.2-005}{633}
对于任意 $n\in\mathbb Z$，令 $L_n$ 为图像 $x=-ny$ 在 $X(B)$ 中的像，沿 $y$ 增加的方向定向。负号来自 Fukaya 范畴的符号约定。由于 $x$ 的周期为一、$y$ 的周期为 $d$，这一图像定义良好，见\BookRef{fig:8.28}{图28}。首先注意，若 $f\colon X(B)\to B$ 为投影，$n_1\neq n_2$，则
\[
f(L_{n_1}\cap L_{n_2})
=f(L_0\cap L_{n_2-n_1})
=B\left(\frac1{n_2-n_1}\mathbb Z\right).
\]
由\BookSectionRef{sec:6.3}{6.3}，期待 $L_n$ 与次数为 $nd$ 的线丛互为镜像；至少除一个符号外如此，符号由 Fukaya 范畴的约定决定，这里并非关键。因此，自然期待 $\widehat{\mathcal L}^{\otimes n}$ 与 $L_n$ 互为镜像。实际上，作为理解这一对偶的第一步，已经在上同调层面，找到两侧间的一个自然同构。
\begin{figure}[H]
\centering
\BookImage{01_Numbered_Figures/Chapter_08/Ch08_Fig28_PDFp634_Bookp624.pdf}
\BookNumberOnlyCaption\label{fig:8.28}
\end{figure}

\SourceParagraph{ch08-8.4.2-006}{634}
更准确地，设 $n_1<n_2$。则
\[
\begin{aligned}
\operatorname{Ext}^i(\widehat{\mathcal L}^{\otimes n_1},
\widehat{\mathcal L}^{\otimes n_2})
&=H^i(\widehat{\mathcal X},\widehat{\mathcal L}^{\otimes(n_2-n_1)})\\
&=\begin{cases}
\displaystyle\bigoplus_{p\in B(\frac1{n_2-n_1}\mathbb Z)}[p]\mathbb C[[q]],&i=0,\\
0,&i=1.
\end{cases}
\end{aligned}
\]
特别地，复形 $\operatorname{Hom}_{D^b_\infty(\widehat{\mathcal X})}
(\widehat{\mathcal L}^{\otimes n_1},\widehat{\mathcal L}^{\otimes n_2})$，拟同构于
\begin{equation}\label{eq:8.15}
\cdots\longrightarrow0\longrightarrow
\bigoplus_{p\in B(\frac1{n_2-n_1}\mathbb Z)}[p]\mathbb C[[q]]
\longrightarrow0\longrightarrow\cdots,
\end{equation}
其中非零项位于次数零。

\SourceParagraph{ch08-8.4.2-007}{634}
另一方面，若 $n_1>n_2$，则
\[
\begin{aligned}
\operatorname{Ext}^i(\widehat{\mathcal L}^{\otimes n_1},
\widehat{\mathcal L}^{\otimes n_2})
&=H^i(\widehat{\mathcal X},\widehat{\mathcal L}^{\otimes(n_2-n_1)})\\
&=\begin{cases}
0,&i=0,\\
\displaystyle\bigoplus_{p\in B(\frac1{n_2-n_1}\mathbb Z)}[p]\mathbb C[[q]],&i=1.
\end{cases}
\end{aligned}
\]
最后一个等式来自 Serre 对偶，后面将给出更显式的描述。因此，态射复形拟同构于仅在次数一非零的复形
\begin{equation}\label{eq:8.16}
\cdots\longrightarrow0\longrightarrow
\bigoplus_{p\in B(\frac1{n_2-n_1}\mathbb Z)}[p]\mathbb C[[q]]
\longrightarrow0\longrightarrow\cdots.
\end{equation}

\SourceParagraph{ch08-8.4.2-008}{634}
与 Fukaya 一侧比较，需要在 $L_n$ 上选择分次。它们在 $X(B)$ 的万有覆盖 $\mathbb R^2$ 中为直线，因此可以取常数分次。按照\BookRef{fig:8.18}{图18}，根据直线的角度，在 $(-1,1)$ 中选择相应分次。于是，由图18可见，$n_1<n_2$ 时，$L_{n_1}\cap L_{n_2}$ 的所有点指标为零；$n_1>n_2$ 时，指标为一。它们当然与 $B(\tfrac1{n_2-n_1}\mathbb Z)$ 一一对应。因此，张量上 $\Lambda_{\mathrm{nov}}$ 后，这两种情形的 Floer 复形，分别与\BookEqRef{eq:8.15}{8.15}、\BookEqRef{eq:8.16}{8.16}一致。

\SourceParagraph{ch08-8.4.2-009}{635}
现在可以比较乘法映射。首先只在 $n_1,n_2>0$ 时，考虑复合
\[
\begin{aligned}
\operatorname{Hom}(\widehat{\mathcal L}^{\otimes(n_0+n_1)},
\widehat{\mathcal L}^{\otimes(n_0+n_1+n_2)})
\otimes\operatorname{Hom}(\widehat{\mathcal L}^{\otimes n_0},
\widehat{\mathcal L}^{\otimes(n_0+n_1)})\\
\longrightarrow
\operatorname{Hom}(\widehat{\mathcal L}^{\otimes n_0},
\widehat{\mathcal L}^{\otimes(n_0+n_1+n_2)}).
\end{aligned}
\]
因为这时可以使用 theta 函数的显式描述。引入 Čech 复形计算高阶乘法映射后，再处理其余情形。

\SourceParagraph{ch08-8.4.2-010}{635}
不妨取 $n_0=0$，要计算的乘积便为
\[
H^0(\widehat{\mathcal L}^{\otimes n_1})
\otimes H^0(\widehat{\mathcal L}^{\otimes n_2})
\longrightarrow H^0(\widehat{\mathcal L}^{\otimes(n_1+n_2)}).
\]
取 $p_1\in B(\tfrac1{n_1}\mathbb Z)$、$p_2\in B(\tfrac1{n_2}\mathbb Z)$，选择其在 $\mathbb R$ 上的提升。暂记
\[
N=n_1+n_2,\qquad
P_\alpha=\frac{n_1p_1+n_2(p_2+\alpha d)}N,
\qquad\alpha=s_2-s_1.
\]
则逐步计算得
\[
\begin{aligned}
\vartheta_{n_1,p_1}\vartheta_{n_2,p_2}
&=\left(\sum_{s_1\in\mathbb Z}z^{T_{n_1}^{s_1}(n_1p_1,n_1\varphi(p_1))}\right)
\left(\sum_{s_2\in\mathbb Z}z^{T_{n_2}^{s_2}(n_2p_2,n_2\varphi(p_2))}\right)\\
&=\sum_{s_1,s_2\in\mathbb Z}
z^{(n_1(p_1+s_1d),n_1\varphi(p_1+s_1d))}
z^{(n_2(p_2+s_2d),n_2\varphi(p_2+s_2d))}\\
&=\sum_{s_1,s_2\in\mathbb Z}
z^{(N(P_\alpha+s_1d),\ n_1\varphi(p_1+s_1d)+n_2\varphi(p_2+s_2d))}\\
&=\sum_{s_1,s_2\in\mathbb Z}
z^{(N(P_\alpha+s_1d),\ N\varphi(P_\alpha+s_1d))}\\
&\qquad\cdot
q^{n_1\varphi(p_1+s_1d)+n_2\varphi(p_2+s_2d)-N\varphi(P_\alpha+s_1d)}.
\end{aligned}
\]

\SourceParagraph{ch08-8.4.2-011}{635}
利用\BookEqRef{eq:8.12}{8.12}，可以验证
\[
\begin{aligned}
&n_1\varphi(p_1+s_1d)+n_2\varphi(p_2+s_2d)
-N\varphi(P_{s_2-s_1}+s_1d)\\
&\qquad=n_1\varphi(p_1)+n_2\varphi(p_2+(s_2-s_1)d)
-N\varphi(P_{s_2-s_1}).
\end{aligned}
\]
对于 $p_1\in\tfrac1{n_1}\mathbb Z$、$p_2\in\tfrac1{n_2}\mathbb Z$，定义
\[
\deg(p_1,p_2)=n_1\varphi(p_1)+n_2\varphi(p_2)
-(n_1+n_2)\varphi\left(\frac{n_1p_1+n_2p_2}{n_1+n_2}\right).
\]

\SourceParagraph{ch08-8.4.2-012}{636}
由 $\varphi$ 的凸性，此数始终非负。因此，
\[
\begin{aligned}
\vartheta_{n_1,p_1}\vartheta_{n_2,p_2}
&=\sum_{\alpha\in\mathbb Z}\sum_{s\in\mathbb Z}
z^{T_N^s(NP_\alpha,N\varphi(P_\alpha))}
q^{\deg(p_1,p_2+\alpha d)}\\
&=\sum_{\alpha\in\mathbb Z}
\vartheta_{N,P_\alpha}q^{\deg(p_1,p_2+\alpha d)}.
\end{aligned}
\]

\SourceParagraph{ch08-8.4.2-013}{636}
几何上，可以如下解释。$p_1,p_2$ 被看作 $B$ 的万有覆盖 $\mathbb R$ 上的点，只定义到模 $d$。因此，对于 $\alpha\in\mathbb Z$，可以取提升 $p_1,p_2+\alpha d$，连接它们的有向区间映到 $B$ 上连接 $p_1,p_2$ 的仿射线段。对所有 $\alpha$ 求和，就相当于对端点为 $p_1\bmod d\mathbb Z$、$p_2\bmod d\mathbb Z$ 的全部仿射线段求和。注意，每条线段的贡献位于一个 $B$ 中的点，它是提升端点的加权平均。

\SourceParagraph{ch08-8.4.2-014}{636}
\BookStatementNumber{18}
\begin{note}[numbered]\label{note:8.18}
这是射影环面簇的一个非常简单现象的推广。若紧格多胞体 $\Delta$ 定义射影环面簇 $(\mathbb P_\Delta,\mathcal O_{\mathbb P_\Delta}(1))$，则 $\Delta(\tfrac1n\mathbb Z)$ 的点，构成 $H^0(\mathbb P_\Delta,\mathcal O_{\mathbb P_\Delta}(n))$ 的一组基。两个基截面 $m_1\in\Delta(\tfrac1{n_1}\mathbb Z)$、$m_2\in\Delta(\tfrac1{n_2}\mathbb Z)$ 的乘积，给出加权平均
\[
\frac{n_1m_1+n_2m_2}{n_1+n_2}
\in\Delta\left(\frac1{n_1+n_2}\mathbb Z\right).
\]
\end{note}

\SourceParagraph{ch08-8.4.2-015}{636}
现在可以与 Fukaya 一侧联系。不妨取 $n_0=0$，考虑乘法映射
\[
\begin{aligned}
\operatorname{Hom}_{\operatorname{Fuk}(\check X)}(L_{n_1},L_{n_1+n_2})
\otimes\operatorname{Hom}_{\operatorname{Fuk}(\check X)}(L_0,L_{n_1})\\
\longrightarrow\operatorname{Hom}_{\operatorname{Fuk}(\check X)}(L_0,L_{n_1+n_2}).
\end{aligned}
\]
如前所述，$p_1,p_2$ 分别确定从 $L_0$ 到 $L_{n_1}$、从 $L_{n_1}$ 到 $L_{n_1+n_2}$ 的态射。给定 $\alpha$，将三条 Lagrange 曲线提升到 $\check X=X(B)=\mathbb C/\langle d,\mathrm i\rangle$ 的万有覆盖 $\mathbb C$，如\BookRef{fig:8.29}{图29}所示。加权平均出现于 $\widetilde L_0\cap\widetilde L_{n_1+n_2}$。因此，阴影三角形 $T$，作为全纯圆盘，对 Floer 乘积的贡献为 $q^{\int_T\omega}$，按照\BookSectionRef{sec:8.3.3}{8.3.3}，其符号为正。图中的三角形面积为
\[
\begin{aligned}
\int_T\omega
&=\frac12\left(\frac{n_1p_1+n_2(p_2+\alpha d)}{n_1+n_2}-p_1\right)
n_1(p_2+\alpha d-p_1)\\
&=\frac12\frac{n_1n_2((p_2+\alpha d)-p_1)^2}{n_1+n_2}\\
&=n_1\psi(p_1)+n_2\psi(p_2+\alpha d)
-(n_1+n_2)\psi\left(\frac{n_1p_1+n_2(p_2+\alpha d)}{n_1+n_2}\right),
\end{aligned}
\]
其中
\begin{equation}\label{eq:8.17}
\psi(x)=\frac{x(x-1)}2.
\end{equation}
这很容易验证。
\begin{figure}[H]
\centering
\BookImage{01_Numbered_Figures/Chapter_08/Ch08_Fig29_PDFp637_Bookp627.pdf}
\BookNumberOnlyCaption\label{fig:8.29}
\end{figure}

\SourceParagraph{ch08-8.4.2-016}{637}
此面积与 $\deg(p_1,p_2+\alpha d)$ 不同，乍看并未得到相同贡献。不过，令
\begin{equation}\label{eq:8.18}
A(x)=\psi(x)-\varphi(x),
\end{equation}
则 $A(x+d)=A(x)$。对于任意 $n_1\neq n_2$，可以在态射的上同调层面定义映射
\[
\begin{aligned}
\operatorname{Ext}^*(\widehat{\mathcal L}^{\otimes n_1},
\widehat{\mathcal L}^{\otimes n_2})\otimes_{\mathbb C[[q]]}\Lambda_{\mathrm{nov}}
\longrightarrow\operatorname{Hom}_{\operatorname{Fuk}(X(B))}(L_{n_1},L_{n_2}),
\end{aligned}
\]
它对 $p\in B(\tfrac1{n_2-n_1}\mathbb Z)$，给出
\begin{equation}\label{eq:8.19}
[p]\longmapsto[p]q^{-(n_2-n_1)A(p)}.
\end{equation}
这就解决了问题。在上述
\[
H^0(\mathcal L^{\otimes n_1})\otimes H^0(\mathcal L^{\otimes n_2})
\longrightarrow H^0(\mathcal L^{\otimes(n_1+n_2)})
\]
的乘积中，$[p_1]$ 与 $[p_2]$ 的乘积在 $[P_\alpha]$ 处的系数为 $q^{\deg(p_1,p_2+\alpha d)}$，其中仍取 $P_\alpha=\tfrac{n_1p_1+n_2(p_2+\alpha d)}{n_1+n_2}$。

\SourceParagraph{ch08-8.4.2-017}{638}
通过\BookEqRef{eq:8.19}{8.19}，将此乘积输送到
\[
\begin{aligned}
\operatorname{Hom}_{\operatorname{Fuk}(X(B))}^0(L_{n_1},L_{n_1+n_2})
\otimes\operatorname{Hom}_{\operatorname{Fuk}(X(B))}^0(L_0,L_{n_1})\\
\longrightarrow\operatorname{Hom}_{\operatorname{Fuk}(X(B))}^0(L_0,L_{n_1+n_2}),
\end{aligned}
\]
则 $[p_2+\alpha d]\otimes1$ 与 $[p_1]\otimes1$ 复合所得 $[P_\alpha]$ 的系数，为
\[
q^{n_1A(p_1)}q^{n_2A(p_2)}q^{\deg(p_1,p_2+\alpha d)}
q^{-(n_1+n_2)A(P_\alpha)}=q^{\operatorname{Area}(T)}.
\]
因此，theta 函数的乘法与 Floer 乘法一致。

\SourceParagraph{ch08-8.4.2-018}{638}
现在希望理解高阶乘积，为此必须具体使用 $D^b_\infty(\widehat{\mathcal X})$ 的 Čech 描述。不过，先给出 Fukaya 一侧的几何描述，使其与上述乘法解释相契合。
